\documentclass[10pt,a4paper]{amsart}
\usepackage{amsmath,amssymb,amsthm,graphicx}
\usepackage{verbatim}
\newenvironment{reference}{\comment}{\endcomment}
%\newenvironment{reference}{}{}
\newenvironment{slogan}{\comment}{\endcomment}
\newenvironment{history}{\comment}{\endcomment}

% For commutative diagrams we use Xy-pic
\usepackage[all]{xy}

% We use 2cell for 2-commutative diagrams.
\xyoption{2cell}
% Theorem environments.
%
\theoremstyle{plain}
\newtheorem{theorem}[subsection]{Theorem}
\newtheorem{proposition}[subsection]{Proposition}
\newtheorem{lemma}[subsection]{Lemma}

\theoremstyle{definition}
\newtheorem{definition}[subsection]{Definition}
\newtheorem{example}[subsection]{Example}
\newtheorem{exercise}[subsection]{Exercise}
\newtheorem{situation}[subsection]{Situation}

\theoremstyle{remark}
\newtheorem{remark}[subsection]{Remark}
\newtheorem{remarks}[subsection]{Remarks}

\numberwithin{equation}{subsection}

% Macros
%
\def\lim{\mathop{\mathrm{lim}}\nolimits}
\def\colim{\mathop{\mathrm{colim}}\nolimits}
\def\Spec{\mathop{\mathrm{Spec}}}
\def\Hom{\mathop{\mathrm{Hom}}\nolimits}
\def\Ext{\mathop{\mathrm{Ext}}\nolimits}
\def\SheafHom{\mathop{\mathcal{H}\!\mathit{om}}\nolimits}
\def\SheafExt{\mathop{\mathcal{E}\!\mathit{xt}}\nolimits}
\def\Sch{\mathit{Sch}}
\def\Mor{\mathop{\mathrm{Mor}}\nolimits}
\def\Ob{\mathop{\mathrm{Ob}}\nolimits}
\def\Sh{\mathop{\mathit{Sh}}\nolimits}
\def\NL{\mathop{N\!L}\nolimits}
\def\CH{\mathop{\mathrm{CH}}\nolimits}
\def\proetale{{pro\text{-}\acute{e}tale}}
\def\etale{{\acute{e}tale}}
\def\QCoh{\mathit{QCoh}}
\def\Ker{\mathop{\mathrm{Ker}}}
\def\Im{\mathop{\mathrm{Im}}}
\def\Coker{\mathop{\mathrm{Coker}}}
\def\Coim{\mathop{\mathrm{Coim}}}

% Boxtimes
%
\DeclareMathSymbol{\boxtimes}{\mathbin}{AMSa}{"02}

%
% Macros for moduli stacks/spaces
%
\def\QCohstack{\mathcal{QC}\!\mathit{oh}}
\def\Cohstack{\mathcal{C}\!\mathit{oh}}
\def\Spacesstack{\mathcal{S}\!\mathit{paces}}
\def\Quotfunctor{\mathrm{Quot}}
\def\Hilbfunctor{\mathrm{Hilb}}
\def\Curvesstack{\mathcal{C}\!\mathit{urves}}
\def\Polarizedstack{\mathcal{P}\!\mathit{olarized}}
\def\Complexesstack{\mathcal{C}\!\mathit{omplexes}}
% \Pic is the operator that assigns to X its picard group, usage \Pic(X)
% \Picardstack_{X/B} denotes the Picard stack of X over B
% \Picardfunctor_{X/B} denotes the Picard functor of X over B
\def\Pic{\mathop{\mathrm{Pic}}\nolimits}
\def\Picardstack{\mathcal{P}\!\mathit{ic}}
\def\Picardfunctor{\mathrm{Pic}}
\def\Deformationcategory{\mathcal{D}\!\mathit{ef}}
\usepackage{cleveref}
\crefname{originaltheorem}{original theorem}{original theorems}
\Crefname{originaltheorem}{Original theorem}{Original theorems}
\makeatletter
\newlabel{theorem-plus}{{00WB}{source}{}{sourcepointer-1}{}}
\newlabel{theorem-plus@cref}{{[originaltheorem][0][]00WB}{[1][0][]source}{}{}{}}
\makeatother
\crefname{originalremark}{original remark}{original remarks}
\Crefname{originalremark}{Original remark}{Original remarks}
\makeatletter
\newlabel{remark-sheaf-condition-empty-covering}{{04B3}{source}{}{sourcepointer-2}{}}
\newlabel{remark-sheaf-condition-empty-covering@cref}{{[originalremark][0][]04B3}{[1][0][]source}{}{}{}}
\makeatother
\crefname{originaldefinition}{original definition}{original definitions}
\Crefname{originaldefinition}{Original definition}{Original definitions}
\makeatletter
\newlabel{definition-site}{{00VH}{source}{}{sourcepointer-3}{}}
\newlabel{definition-site@cref}{{[originaldefinition][0][]00VH}{[1][0][]source}{}{}{}}
\makeatother
\title{Characterizations of sheaf theoretically empty objects}
\author{The Stacks Project Authors}
\date{stacks-project:sites.tex:local-snapshot-20261012; GNU FDL 1.2 or later}
\begin{document}
\maketitle
\noindent\textit{Editorial note.} All delivered mathematical text is the author's own, copied from the pinned Stacks Project chapter file sites-cohomology.tex; no statement or proof step was rewritten, reordered or removed.\par\medskip
\noindent\textit{Editorial note.} The proof is the author's complete argument; every step the author writes out is inside the counted span and no closing sentence is a placeholder.\par\medskip
\noindent\textit{Editorial note.} Every labelled reference occurring inside the counted spans is carried in assembly.external\_references as a non-load-bearing navigation pointer (same-chapter ones at their definition site, cross-chapter ones at the invocation inside the counted span), so no load-bearing mathematics is replaced by a stub. No bibliography citation occurs inside the counted spans.\par\medskip
\noindent\textit{Editorial note.} The checked chapter list for same-name collisions is the set of stacks chapters staged in this repository: chow, cohomology, divisors, etale-cohomology, flat, more-algebra, more-morphisms, morphisms, spaces-more-morphisms, varieties together with this card's chapter sites-cohomology.\par\medskip
\noindent\textit{Editorial note.} The preamble is the chapter's own retained minimal preamble (sources/stacks-sites-cohomology/retained-preamble.tex), a verbatim excerpt of the native preamble: the Xy-pic setup and the native theorem declarations with the mathematical macros.\par\medskip
\noindent\textit{Allowed standard prerequisites (editorial).}\begin{itemize}
\item The plus construction and its relation to sheafification. Same chapter, Theorem theorem-plus. Shows that the plus construction applied to the empty presheaf is its sheafification.
\item The sheaf condition for the empty covering. Same chapter, Remark remark-sheaf-condition-empty-covering. Supplies the implication from the empty family being a covering to values being singletons.
\item The axioms of a site. Same chapter, Definition definition-site. Supplies the stability axiom used for the final assertion.
\end{itemize}\medskip
\ifdefined\hypertarget\hypertarget{sourcepointer-1}{}\fi
\noindent\textit{Source reference.} Original theorem 00WB: sources/stacks-sites/pinned-original-source.tex, source lines 10158--10158. Navigation pointer invoked at proof line 10158; the counted span contains the pointer only and the cited result is not reproduced inside it, so the pointer is registered as a non-load-bearing navigation stub in the established form for Stacks material. No bibliography citation is involved.\par\medskip
\ifdefined\hypertarget\hypertarget{sourcepointer-2}{}\fi
\noindent\textit{Source reference.} Original remark 04B3: sources/stacks-sites/pinned-original-source.tex, source lines 10162--10162. Navigation pointer invoked at proof line 10162; the counted span contains the pointer only and the cited result is not reproduced inside it, so the pointer is registered as a non-load-bearing navigation stub in the established form for Stacks material. No bibliography citation is involved.\par\medskip
\ifdefined\hypertarget\hypertarget{sourcepointer-3}{}\fi
\noindent\textit{Source reference.} Original definition 00VH: sources/stacks-etale-cohomology/pinned-original-source.tex, source lines 526--526. Navigation pointer invoked at proof line 10165; the counted span contains the pointer only and the cited result is not reproduced inside it, so the pointer is registered as a non-load-bearing navigation stub in the established form for Stacks material. No bibliography citation is involved.\par\medskip
\begin{lemma}
\label{lemma-characterize-empty}
Let $\mathcal{C}$ be a site. Let $U$ be an object of $\mathcal{C}$.
The following are equivalent:
\begin{enumerate}
\item $U$ is sheaf theoretically empty,
\item $\mathcal{F}(U)$ is a singleton for each sheaf $\mathcal{F}$,
\item $\emptyset^\#(U)$ is a singleton,
\item $\emptyset^\#(U)$ is nonempty, and
\item the empty family is a covering of $U$ in $\mathcal{C}$.
\end{enumerate}
Moreover, if $U$ is sheaf theoretically empty, then for any morphism
$U' \to U$ of $\mathcal{C}$ the object $U'$ is sheaf theoretically empty.
\end{lemma}

\begin{proof}
For any sheaf $\mathcal{F}$ we have
$\mathcal{F}(U) = \Mor_{\Sh(\mathcal{C})}(h_U^\#, \mathcal{F})$.
Hence, we see that (1) and (2) are equivalent.
It is clear that (2) implies (3) implies (4).
If every covering of $U$ is given by a nonempty family,
then $\emptyset^+(U)$ is empty by definition of the plus construction.
Note that $\emptyset^+ = \emptyset^\#$ as $\emptyset$ is a separated
presheaf, see
Theorem \ref{theorem-plus}.
Thus we see that (4) implies (5). If (5) holds, then
$\mathcal{F}(U)$ is a singleton for every sheaf $\mathcal{F}$
by the sheaf condition for $\mathcal{F}$, see
Remark \ref{remark-sheaf-condition-empty-covering}.
Thus (5) implies (2) and (1) -- (5) are equivalent. The final
assertion of the lemma follows from Axiom (3) of
Definition \ref{definition-site}
applied the empty covering of $U$.
\end{proof}

\end{document}
